\documentclass{article}

% \usepackage{corl_2026} % Use this for the initial submission.
% \usepackage[final]{corl_2026} % Uncomment for the camera-ready ``final'' version.
\usepackage[preprint]{corl_2026} % Uncomment for pre-prints (e.g., arxiv); This is like ``final'', but will remove the CORL footnote.

\usepackage{amsmath,amssymb,amsfonts}
\usepackage{algorithmic}
\usepackage{algorithm}
\usepackage{graphicx}
\usepackage{textcomp}
\usepackage{booktabs}
\usepackage{multirow}
\usepackage{xcolor}
\usepackage{subcaption}

% Custom commands
\newcommand{\methodname}{DiscreteRTC}
\newcommand{\ie}{\textit{i.e.}}
\newcommand{\eg}{\textit{e.g.}}
\newcommand{\etal}{\textit{et al.}}
\newcommand{\todo}[1]{\textcolor{red}{[TODO: #1]}}
\newcommand{\chen}[1]{\textcolor{blue}{[Chen: #1]}}

\title{DiscreteRTC: Discrete Diffusion Policies \\ are Natural Asynchronous Executors}

% NOTE: authors will be visible only in the camera-ready and preprint versions (i.e., when using the option 'final' or 'preprint').
% 	For the initial submission the authors will be anonymized.

\author{
  Pengcheng Wang$^1$, Kaiwen Hong$^2$, Chensheng Peng$^1$,\\ {\bf Katherine Driggs-Campbell}$^2$, {\bf Masayoshi Tomizuka}$^1$, {\bf Chenfeng Xu}$^3$, {\bf Chen Tang}$^4$\\
  UC Berkeley$^1$, UIUC$^2$, UT Austin$^3$, UCLA$^4$\\
  \texttt{wangpc@berkeley.edu, cxu@utexas.edu, ctangac@ucla.edu}
}


\begin{document}
\maketitle
\vspace{-2em}
%===============================================================================

\begin{abstract}
% [SENT] latency problem
% [SENT] async is structural, not optional
Unlike chatbots, physical AI must act while the world keeps evolving.
Therefore, the inter-chunk pause of synchronous executors are fatal for dynamic tasks regardless of how fast the inference is.
% [SENT] async + RTC
Asynchronous execution---\emph{thinking while acting}---is therefore a structural requirement, and real-time chunking (RTC) makes it viable by recasting chunk transitions as inpainting: freezing committed actions and consistently generating the remainder.
% [SENT] flow-matching RTC is structurally suboptimal
However, RTC with flow-matching policy is structurally suboptimal: its inpainting comes from inference-time corrections rather than the base policy, yielding little pre-training benefit, specific fine-tuning, heuristic guidance weights, and extra computation that inflates the latency.
% [SENT] discrete diffusion is natural
In this work, we observe that discrete diffusion policies, which generate actions by iteratively unmasking, are \emph{natural asynchronous executors that resolve all limitations at once}: they are fine-tuning free since inpainting is their native operation, while early stopping further provides adaptive guidance and reduces inference cost.
% [SENT] DiscreteRTC + results
We propose \textbf{DiscreteRTC}, which replaces external corrections with native unmasking, and show on dynamic simulated benchmarks and real-world dynamic manipulation tasks that it achieves higher success rates than continuous RTC and other baselines, while being simpler to implement, faster at inference, and better at execution.
More visualizations are on our \href{https://outsider86.github.io/DiscreteRTCSite/}{website}.
\end{abstract}


\begin{figure}[!ht]
    \centering
    \includegraphics[width=\linewidth]{figures/PnPcover3.png}
    \caption{\textbf{RTC with discrete diffusion policies solving dynamic manipulation.}}
    \label{fig:cover}
    \vskip -0.2in
\end{figure}



%===============================================================================
\section{Introduction}
\label{sec:intro}

% % [PARA 1] Physical AI is advancing rapidly, and real-time performance is non-negotiable.
% % [SENT] AI systems increasingly interact with the physical world
% Physical AI systems have grown more and more capable~\citep{black2024pi_0, intelligence2025pi_, bjorck2025gr00t, wang2026dadp}: from dexterous manipulators~\citep{openai2018learning, an2025dexterous} assembling intricate parts to mobile robots capable of different athletic events~\citep{su2025hitter, zhang2026learningathletichumanoidtennis, yu2025skillmimic}.
% However, this growing agility and intelligence come at the cost of ever-larger model sizes and inference times~\citep{kaplan2020scaling}.
% % [SENT] real-time constraint is fundamental
% Unlike chatbots~\citep{brown2020language, yang2025qwen3} or image generators~\citep{wan2025wan}, these cyber-physical systems always operate in real time: the world keeps evolving while the robot is ``thinking''.
% % [SENT] consequence of delay
% Therefore, delays between sensing and acting have a tangible impact on performance---a slow and laggy policy does not merely annoy a user; it can sometimes cause catastrophic consequences.



% [PARA 1] Async execution is structurally necessary; RTC is the forced consequence.
% [SENT] physical AI operates under a temporal asymmetry
Unlike chatbots~\citep{brown2020language, yang2025qwen3} or image generators~\citep{wan2025wan}, physical AI systems~\citep{black2024pi_0, intelligence2025pi_, bjorck2025gr00t, wang2026dadp, openai2018learning, an2025dexterous, su2025hitter, yu2025skillmimic} act while the world keeps evolving: the robot cannot afford to pause while thinking.
% [SENT] synchronous execution is structurally broken, not just slow
This reveals a structural limitation of synchronous execution with action-chunking policies~\citep{chi2025diffusion, lai2025action}---every inference call forces a mandatory pause between chunks during which the robot is idle but the world is not, which is fatal for dynamic tasks such as tracking moving targets~\citep{xie2026dynamicvla}.
% [SENT] async + chunking as the structural answer
Therefore, asynchronous execution---\emph{thinking while acting} over action chunks---is not a latency optimization but a structural requirement.


% [SENT] naive async breaks + RTC as the principled fix
However, naive async strategies, whether switching chunks immediately or smoothing them via temporal ensembling~\citep{zhao2023learning}, introduce inter-chunk discontinuity or invalid actions because each new chunk is generated in isolation from the committed actions.
Real-time chunking (RTC)~\citep{black2025real} resolves this by recasting chunk transitions as an \emph{inpainting} problem: freezing committed actions from the previous chunk and consistently generating the remainder, representing a simple yet effective approach that has achieved great success in precise and smooth real-world manipulation~\citep{intelligence2025pi_06, intelligence2026pi}.

\newpage
% [PARA 3] RTC with flow-matching is suboptimal; we can do better.
% [SENT] RTC with flow-matching is limited
Specifically, RTC is specially designed for flow-matching policies~\citep{lipman2022flow}.
As the predominant action head architecture in today's state-of-the-art VLAs~\citep{black2024pi_0, intelligence2025pi_, bjorck2025gr00t}, flow-matching has driven much of the recent progress in large-scale robot learning.
However, in this paper, we first systematically show that \emph{\textbf{RTC with a flow-matching head is far from ideal with four critical limitations}}:


% [SENT] (a) pre-training gap
\begin{itemize}
    \item [(a)] \textbf{Pre-training without Inpainting.} Flow-matching policies are not pre-trained on action chunks with inconsistent noise levels for inference-time inpainting.
    Therefore, scaling pre-training does not directly improve asynchronous performance, and inference-time corrections like $\Pi$GDM~\citep{song2023pseudoinverse} are inevitable for effective inpainting;
    \item [(b)] \textbf{Fine-tuning Required.} As a consequence, adequate inpainting quality demands a dedicated fine-tuning stage with techniques such as action-suffix conditioning~\citep{black2025training, wang2026real} to explicitly introduce the inpainting-specific noise pattern into training;
    \item [(c)] \textbf{Heuristic Guidance.} Moreover, to better leverage the previous action chunks, $\Pi$GDM requires a heuristic schedule fixed across different inference-time cases;
    \item [(d)] \textbf{Extra Inference Cost.} Finally, the correction guidance term at every denoising step roughly doubles inference cost, ironically increasing the very latency RTC aims to hide.
\end{itemize}

Unlike many ongoing efforts~\citep{liu2026learning, lu2026faster, black2025training, yang2026abpolicy}, this work is not motivated to seek innovation in solving these issues under the flow-matching paradigm.
Instead, our key observation is that by replacing the action head with a discrete diffusion policy~\citep{liang2025discrete}, all the aforementioned limitations can be resolved at once. Or, to put it simply: \emph{\textbf{Discrete Diffusion Policies are Natural Asynchronous Executors}}:

\begin{itemize}
    \item [(a)] \textbf{Inpainting as Pre-training.} Discrete diffusion policies are pre-trained to inpaint upon randomly masked sequences. Therefore, scaling pre-training directly improves asynchronous performance, and the native forward pass suits inference-time inpainting;
    \item [(b)] \textbf{Fine-tuning Free.} As a consequence, inpainting-specific patterns are implicitly introduced during pre-training, making discrete diffusion a fine-tuning-free approach for high-quality, out-of-the-box asynchronous execution;
    \item [(c)] \textbf{Natural Guidance.} Moreover, with discrete diffusion policies, we can early-exit inference once the necessary action tokens are unmasked, leaving the remaining masking pattern as an adaptive and natural guidance for the next inference;
    \item [(d)] \textbf{Lower Inference Cost.} Finally, with committed tokens from previous chunks, the tokens to unmask per inference is reduced, leading to lower inference cost for inpainting.
\end{itemize}

% [PARA 5] We propose DiscreteRTC and validate it.
% [SENT] we propose DiscreteRTC
We propose the resulting async execution method, \textbf{DiscreteRTC}, which replaces the online correction in flow-matching-based RTC with the native unmasking operation of discrete diffusion policies.
% [SENT] experimental validation
Experiments on the dynamic simulated benchmark Kinetix~\citep{matthews2024kinetix} and real-world dynamic manipulation demonstrate that DiscreteRTC achieves higher success rates than flow-matching-based RTC and other baselines, being simpler to implement, faster to inference, and better to execute.


% [PARA] Summary of contributions.
% [SENT] contributions list
% In summary, our contributions are:
% \begin{enumerate}
%     \item We provide a systematic analysis of the fundamental limitations of the currently popular RTC framework with a flow-matching action head.
%     \item We show that these limitations can be naturally resolved by discrete diffusion policies. In this sense, \emph{\textbf{Discrete Diffusion Policies are Natural Asynchronous Executors}}.
%     \item Through simulation and real-world experiments, we demonstrate that the resulting method, \textbf{DiscreteRTC}, (a) achieves higher task success rates and throughputs under varying delays, (b) outperforms training-time methods that require fine-tuning, (c) requires less inference time, and (d) can be integrated with advanced techniques such as VLASH~\citep{tang2025vlash}.
% \end{enumerate}

%===============================================================================
\section{Preliminaries}
\label{sec:background}

\subsection{Real-time Chunking}
\label{sec:bg_rtc}

% [PARA] Action chunking outputs a sequence of actions per inference call.
% [SENT] definition of action chunking
\textbf{Action Chunking Flow Policy.}
We consider an action chunking policy $\pi$ that receives an observation $\mathbf{o}_t$ and outputs a chunk of $H$ future actions $\mathbf{A}_t = (a_t, a_{t+1}, \ldots, a_{t+H-1})$.
State-of-the-art VLAs~\citep{black2024pi_0, intelligence2025pi_} generate action chunks using conditional flow matching~\citep{lipman2022flow, zhan2026mean}: noise $\mathbf{A}^0_t \sim \mathcal{N}(0, I)$ is sampled and integrated through the learned velocity field $v_\pi$ from $\tau=0$ to $1$:
\begin{equation}
    \mathbf{A}_t^{\tau+\frac{1}{n}} = \mathbf{A}_t^\tau + \frac{1}{n} v_\pi( \mathbf{A}_t^\tau, \mathbf{o}_t, \tau),
    \label{eq:flow_step}
\end{equation}
where $\tau \in [0, 1)$ is the flow timestep and $n$ is the number of denoising steps.
% [SENT] training objective
During training, all actions within a chunk are corrupted at the same noise level $\tau$ and denoised together.

% [PARA] RTC runs inference asynchronously and uses inpainting for continuity.
% [SENT] motivation for async execution
\textbf{Real-Time Chunking.}
Running modern VLAs with action generation time $\delta$ smaller than the controller period $\Delta t$ (typically 20$\sim$50\,ms) is nearly impossible~\citep{black2025real}: even a 7B OpenVLA model optimized for speed~\citep{kim2025fine, kim2024openvla} requires $\delta = 321$\,ms on a server-grade A100 GPU.
% [SENT] async execution and discontinuity
The robot must therefore act while the policy is still computing, \ie, asynchronous execution.
Naively switching to a new action chunk as soon as it is ready, however, leads to inter-chunk discontinuity: adjacent chunks may correspond to different behavioral modes, producing jerky motions at chunk boundaries.

% [SENT] RTC overview
Real-time chunking (RTC)~\citep{black2025real} addresses this discontinuity by posing asynchronous execution as an \emph{inpainting} problem.
% [SENT] delay and frozen actions
If inference takes $d$ (inference delay) steps to complete, the first $d$ actions of the new chunk are ``frozen'' to the values from the previous chunk before the new chunk becomes available.
% [SENT] inpainting formulation
The remaining actions must be generated to be consistent with this frozen prefix.
% [SENT] ΠGDM guidance steers generation toward the frozen prefix
For flow-matching policies, RTC implements this inpainting via $\Pi$GDM~\citep{song2023pseudoinverse}, adding a gradient-based correction at each denoising step that steers the generation toward the target prefix $\mathbf{Y}$:
\begin{equation}
    \mathbf{g} = \left(\mathbf{Y} - \hat{f}_{\mathbf{A}^1}(\mathbf{A}^\tau)\right)^\top \text{diag}(\mathbf{W}) \cdot \frac{\partial \hat{f}_{\mathbf{A}^1}}{\partial \mathbf{A}'}\bigg|_{\mathbf{A}^\tau},
    \label{eq:guidance}
\end{equation}
where $\hat{f}_{\mathbf{A}^1}$ denotes the one-step denoising function and $\mathbf{W}$ is a vector of soft-masking weights.
% [SENT] soft masking with exponential decay
The weights $\mathbf{W}$ assign $1$ to the first $d$ frozen actions, exponentially decay for the overlapping region, and $0$ for the final $s$ (execution horizon) actions.
During inference, the robot starts an inference every $s$ steps and execute the next $s$ actions from the most recent inference results.


\subsection{Discrete Diffusion Policies}
\label{sec:bg_discrete}

% [PARA] Discrete diffusion generates sequences by iteratively unmasking tokens.
% [SENT] tokenization of continuous actions
Discrete diffusion policies~\citep{liang2025discrete} first tokenize continuous actions into discrete tokens via a learned or fixed vocabulary (\eg, VQ-VAE~\citep{VQVAE}, FAST~\citep{FAST} or $k$-bin quantization).
With a slight abuse of notation, we reuse $\mathbf{A}_t$ to denote the resulting action token sequence.
% [SENT] forward process masks tokens
The forward process replaces tokens with a special $[\text{MASK}]$ token with a masking schedule,
% [SENT] reverse process predicts masked tokens conditioned on unmasked ones
while the reverse process iteratively executes unmasking operations conditioned on the current partially masked token sequence.

% [SENT] two components
A discrete diffusion policy is trained for one-step unmasking operation.
It consists of two components: $p_\pi$, which predicts the probability distribution over all token positions, and $f_\pi$, which samples which tokens to unmask at the current step from that distribution.
At each step $k$, the policy first predicts token identities via $p_\pi$ and then selects which tokens to unmask via $f_\pi$:
\begin{equation}
    \mathbf{A}_t^{k+1} = \pi(\mathbf{A}_t^k, \mathbf{o}_t) = f_\pi\!\bigl(p_\pi(\mathbf{A}_t^k, \mathbf{o}_t)\bigr),
    \label{eq:dd_step}
\end{equation}
where $\mathbf{A}_t^k$ is the partially unmasked token sequence at step $k$ and $\mathbf{o}_t$ is the observation.
% [SENT] training objective
During training, the discrete diffusion policy learns to predict the fully unmasked ground-truth action token sequence given the randomly masked token sequence.


%===============================================================================

\section{Flow-Matching is not Suitable for RTC}
\label{sec:limitations}

\begin{figure}
    \centering
    \includegraphics[width=\linewidth]{figures/continuousRTC.pdf}
    \caption{\textbf{RTC with flow-matching head.} Color represents the noise level, where green stands for the clear action and yellow stands for the pure noise. The flow-matching head is ill-suited for RTC because (a) during pre-training, the base policy is not trained on inpainting tasks; (b) to acquire this capability, a specially designed fine-tuning stage is required; (c) at inference time, RTC relies on heuristic guidance weight design; and (d) incurs extra cost for online correction computation.}
    \label{fig:continuousRTC}
    \vskip -0.15in
\end{figure}

% [PARA] Overview: four limitations of applying flow-matching to RTC.
% [SENT] section roadmap
As illustrated in Figure~\ref{fig:continuousRTC}, we analyze four critical limitations of RTC with flow-matching policies.

%\newpage
\subsection{Pre-training without Inpainting}
\label{sec:lim_pretrain}
% [PARA] Flow-matching pre-training assumes uniform noise across the chunk, which conflicts with the mixed-noise pattern required by RTC.
% [SENT] standard training uses uniform noise level
Pre-training of flow-matching policies~\citep{lipman2022flow} learns to generate action chunks conditioned on observations and a noised action chunk.
% [SENT] all actions share the same noise level during training
During training, all actions within a chunk are corrupted at the \textbf{consistent} noise level $\tau$ and are denoised together under the same computational budget.
% [SENT] RTC requires mixed noise levels across the chunk
However, during RTC, the generation process starts from an action chunk with \textbf{inconsistent} noise levels across positions: 0 for the committed actions, 1 for the new actions, and interpolated values for the intermediate actions.
% [SENT] guidance is external and unseen during pre-training
Although the $\Pi$GDM~\citep{song2023pseudoinverse} guidance provides an external correction at inference time for effective inpainting, the model has encountered neither such corrections nor the mixed-noise chunk pattern during the pre-training stage.
% [SENT] scaling pre-training does not help inpainting
As a result, scaling up pre-training data or model capacity does not directly improve inpainting quality for asynchronous execution.

\newpage
\subsection{Fine-tuning Required}
\label{sec:lim_finetuning}
% [PARA] Closing the pre-training gap requires a dedicated fine-tuning stage with mixed-noise chunks.
% [SENT] mixed-noise chunks must be introduced during training
Following Sec.~\ref{sec:lim_pretrain}, to improve the inpainting capability of the base policy, the mixed-noise chunk pattern must be explicitly introduced during training along with a corresponding inpainting objective.
% [SENT] fine-tuning adds complexity and may hurt base quality
This requires fine-tuning the model with an explicit inpainting loss~\citep{black2025training}, adding training complexity and risking interference with the base generation quality~\citep{wang2026real}.
% [SENT] engineering and compute burden for large VLAs
This introduces a non-trivial additional engineering and computational burden, particularly for large VLAs that are already expensive to fine-tune.
% [SENT] large-scale pre-training may demand large-scale fine-tuning
Moreover, the fine-tuning recipe itself must be carefully designed~\citep{hu2022lora}: for models pre-trained at large scale, a correspondingly large-scale fine-tuning effort may be necessary to reliably acquire the inpainting capability without degrading other abilities~\citep{wang2026real}.

\subsection{Heuristic Guidance}
\label{sec:lim_guidance}

To effectively inpaint with flow-matching policies, online corrections like $\Pi$GDM~\citep{song2023pseudoinverse} are inevitable.
Completing Equation~\ref{eq:guidance}, the full denoising step of $\Pi$GDM is presented as follows:

\begin{equation}
\label{eqn: continuousRTC forward}
\mathbf{A}_t^{\tau+\frac{1}{n}} = \mathbf{A}_t^\tau + \frac{1}{n}\!\left(v_\pi(\mathbf{A}_t^\tau, \mathbf{o}_t, \tau) + \min\!\left(\beta,\, \frac{1-\tau}{\tau \cdot r_\tau^2}\right) \left(\mathbf{Y} - \hat{f}_{\mathbf{A}^1}(\mathbf{A}^\tau)\right)^\top \text{diag}(\mathbf{W}) \frac{\partial \hat{f}_{\mathbf{A}^1}}{\partial \mathbf{A}'}\bigg|_{\mathbf{A}^\tau} \right).
\end{equation}
% [PARA] The guidance weights and clipping threshold are hand-designed heuristics that do not adapt to the inference situation.
% [SENT] soft masking uses a hand-designed exponential decay
Notably, the soft-masking weights $W$ in Equation~\ref{eqn: continuousRTC forward} follow a hand-crafted exponential decay schedule.
% [SENT] the schedule is fixed and only empirically validated
This important guidance schedule, which decides how to follow the previous action chunks, is fixed across different inference situations (e.g., varying delays, execution horizons, and tasks) and is validated only through empirical studies.
% [SENT] tuning burden across configurations
This heuristic design not only creates a tuning burden, but also behaves inflexibly in response to different inference cases during asynchronous execution.

\subsection{Extra Inference Cost}
\label{sec:lim_cost}
% [PARA] Gradient-based guidance roughly doubles per-step cost, ironically increasing the very latency it aims to hide.
% [SENT] VJP through the denoising network at every step
Finally, running $\Pi$GDM at inference time requires non-negligible extra computation.
At each denoising step, computing the guidance term in Equation~\ref{eqn: continuousRTC forward} requires a vector-Jacobian product (VJP), which roughly doubles the per-step computation~\citep{wang2026real}. As a mechanism intended to enable smooth real-time execution, flow-matching-based RTC instead introduces additional inference delay, shrinking the time window available for generating new actions and counteracting its design goal.

%===============================================================================

\section{Discrete Diffusion Policies are Natural Asynchronous Executors}
\label{sec:method}

\begin{figure}
    \centering
    \includegraphics[width=\linewidth]{figures/discreteRTC.pdf}
    \caption{\textbf{RTC with discrete diffusion head.}  Color represents the masking status, where green stands for the unmasked token and yellow stands for the masked token. The discrete diffusion head is naturally suited for RTC because (a) during pre-training, the base policy is already trained on inpainting tasks; (b) consequently, no inpainting-specific fine-tuning is required; (c) at inference time, early stopping from the previous inference leaves a natural guidance signal; and (d) fewer tokens need to be unmasked per step, reducing the computational cost of each forward pass.}
    \label{fig:discreteRTC}
    \vskip -0.15in
\end{figure}

% [PARA] Section overview: discrete diffusion naturally resolves the limitations of flow-matching RTC.
% [SENT] section roadmap
As illustrated in Figure~\ref{fig:discreteRTC}, we show that discrete diffusion policies can naturally resolve all the aforementioned limitations of flow-matching-based RTC at once.
\subsection{Inpainting as Pre-training \& Fine-tuning Free}
\label{sec:method_pretrain}

% [PARA] Inpainting is the native operation of discrete diffusion, resolving (a) and (b).
% [SENT] inpainting is structurally identical to standard masked generation
In discrete diffusion, inpainting is the native operation: given a partially masked token sequence, the policy is asked to reconstruct the target action chunks by predicting the masked tokens.
% [SENT] resolves (a): pre-training directly improves inpainting
Therefore, pre-training directly improves the model's ability to inpaint at inference time: scaling model, data, and pre-training conputation yield better asynchronous execution without inpainting-specific fine-tuning.
% [SENT] no modification needed for inpainting
To inpaint at inference time, no extra implementation, loss term, or inference-time guidance are needed. Unlike Equation~\ref{eqn: continuousRTC forward}, the model generates conditioned on partially unmasked token sequence from last inference, exactly as it was trained --- simply forwarding the base policy:
\begin{equation}
    \label{eq:discrete_rtc_forward}
    \mathbf{A}_t^{k+1} = \pi(\mathbf{A}_t^k, \mathbf{o}_t).
\end{equation}



\subsection{Natural Guidance}
\label{sec:method_netural}



% [PARA] Unmasked tokens carry immediate semantic meaning, enabling early stopping.
% [SENT] key property: unmasked tokens are already valid
To solve the guidance issue, we exploit a key property of discrete diffusion: once a token is unmasked, it already corresponds to a clear action semantic.
% [SENT] contrast with continuous noisy actions
This stands in contrast to continuous diffusion or flow-matching, where the denoised action chunk becomes valid and executable all together only when the noise level reaches zero at the end of the denoising process. 
Therefore, with discrete diffusion, the model can not only start inference from an intermediate state (e.g. a partially unmasked action chunk with prefix), but also terminate at an intermediate stage. 

During async execution, to ensure valid execution before the next inference, only the next $s$ actions after the first $d$ commited actions are necessary to be fully unmasked. Therefore, the inference can early-exit when satisfying the minimal criteria, leaving the remaining $H - d - s$ actions a partially unmasked pattern. By directly carrying over the results, this pattern serves as a \emph{natural and adaptive} guidance signal for the next inference with the number of the unmasked tokens, replacing the heuristic and fixed guidance schedule in flow-matching basd RTC.
% [SENT] both intermediate start and early stop are possible

% [PARA] Early stopping provides natural guidance and reduces cost.
% [SENT] resolves (c): partially unmasked actions serve as natural adaptive guidance

% [SENT] remaining tokens carry forward as guidance for next inference


% [PARA] Resolves (d): inference cost is reduced rather than increased.
% [SENT] fewer unmasking steps needed
\subsection{Lower Inference Cost}
\label{sec:method_lower}
Finally, since unmasking starts from an intermediate state and ends at an intermediate stage, keeping the number of tokens unmasked per step constant actually reduces the total number of tokens to unmask per inference to roughly $s / H$ of the original, or at least $1 - d/H$ by starting only from the actions to execute and terminating at complete action chunks.
% [SENT] cost reduction instead of cost increase
This means that, instead of \emph{increasing} inference cost as in flow-matching-based RTC, our method can actually \emph{reduce} it by simply forwardin the base policy with fewer interative steps.
% [SENT] alternative: finer-grained actions at same cost
Alternatively, we can keep the number of inference steps fixed and reduce the number of tokens unmasked per step, thereby producing finer-grained actions without increasing the action generation cost compared with generation from scratch.

%===============================================================================
\begin{figure}
    \centering
    \includegraphics[width=\linewidth]{figures/KinetixMain.pdf}
    \caption{\textbf{Experimental Results in Kinetix.} The throughputs represent the task completed by the policy every 256 steps.
    \textbf{Left}: Average solve rate and throughputs across all environments with different inference delays;
    \textbf{Right}: Solve rates for every tasks with different inference delays.
    The executions horizon follows $s = \max(1, d)$ and each datapoint represents 2048 trials.}
    \label{fig:kinetix_main}
    \vskip -0.15in
\end{figure}

\section{Experiments}
\label{sec:experiments}




% % [PARA] Experimental questions framing the evaluation.
% % [SENT] three experimental questions
% In our experiments, we aim to answer the following questions:

% \begin{itemize}
%     \item How does DiscreteRTC perform comparing to existing inference-time and training-time asynchronous execution methods under increasing inference delays?
%     \item How does the DiscreteRTC effect the performance and efficiency both in simulation and real-world dynamic tasks?
% \end{itemize}

\subsection{Simulated Benchmark}
\label{sec:exp_sim}

% [PARA] Benchmark design: dynamic tasks where asynchronous execution is necessary.
% [SENT] quasi-static benchmarks are insufficient
% [SENT] we use dynamic tasks with force-based control
We first evaluate DiscreteRTC in Kinetix~\citep{matthews2024kinetix}, a simulated benchmark of dynamic tasks.
% [SENT] tasks involve dynamic motions
All tasks involve dynamic motions such as throwing, catching, and balancing.
% [SENT] Gaussian noise makes closed-loop corrections critical
To simulate imperfect actuation, we add Gaussian noise to the actions to make closed-loop corrections crucial for success~\citep{wang2026real}. Please refer to Appendix~\ref{app: addexp} for more results.

% [PARA] Data generation and policy training.
% [SENT] expert policies via RL
\textbf{Setup.}
We implemented the discrete diffusion policies upon the official codebase and datasets from RTC~\citep{black2025real}.
% [SENT] dataset generation from multiple experts
For fair comparison, we adopt the same training setup (optimizer, learning rate, epochs, etc.) and the same policy network size except the final layer for logits output.
% [SENT] discrete diffusion policy training
We use a trivial $512$-bin quantization for the action chunk discretization.
Please refer to Appendix~\ref{app: simexp} for more details.

% [PARA] Baselines covering the main alternative approaches.
% [SENT] baseline descriptions
\textbf{Baselines.} With same policy architecture, we compare against the flow-matching (Continuous) and the discrete diffusion (Discrete) action head under the following asynchronous execution strategies:
\begin{itemize}

    \item \textbf{Naive Async.} Chunks are generated from scratch and switched when the new one is ready.
    \item \textbf{Bidirectional decoding (BID)}~\citep{liu2024bidirectional}. Uses rejection sampling to maintain cross-chunk continuity. We use a batch size of $N = 16$ with no weak policy as the official implementation.
    \item \textbf{RTC}~\citep{black2025real}. Flow-matching policy with $\Pi$GDM inpainting and soft masking, representing the state-of-the-art asynchronous execution method for continuous policies.
\end{itemize}





% \begin{figure}[t]
%     \centering
%     \begin{minipage}{0.5\textwidth}
%         \centering
%         \includegraphics[width=0.9\linewidth]{figures/KinetixStep.pdf}
%         \caption{Required iterative steps of different policy architectures in Kinetix with $s = \max(1, d)$.}
%         \label{fig: kinetix_steps}
%     \end{minipage}
%     \hfill
%     \begin{minipage}{0.45\textwidth}
%         \centering
%         \includegraphics[width=\linewidth]{figures/KinetixAblation.pdf}
%         \caption{Experimental results of extended variants in Kinetix with $s = \max(1, d)$.}
%         \label{fig: kinetix_abl}
%     \end{minipage}
%     \vskip -0.15in
% \end{figure}


\begin{figure}
    \centering
    \includegraphics[width=\linewidth]{figures/Step&Ablation.pdf}
    \caption{\textbf{Extend Experimental Results in Kinetix.} The evaluation setup keeps the same with Figure~\ref{fig:kinetix_main}.
    \textbf{Left}: Required iterative steps of different poliicy architectures in Kinetix with $s = \max(1, d)$;
    \textbf{Right}: Average solve rates of extended variants in Kinetix.
    }
    \label{fig:step_ablation}
    \vskip -0.2in
\end{figure}

% [PARA] Main simulation results.
% [SENT] results figure reference

\textbf{Results.}
We present the main results in Kinetix in Figure~\ref{fig:kinetix_main}.
% [SENT] discrete policies outperform continuous counterparts overall
Under different inference delays, DiscreteRTC consistently outperforms ContinuousRTC and other variants on both solve rates and throughputs, showing the advantage of the native inpainting capability in discrete diffusion policies.



% [SENT] ablation and compatibility with other methods
We present the extended experimental results in Kinetix in Figure~\ref{fig:step_ablation}.
As shown in the results of iterative steps per policy, DiscreteRTC requires fewer iterative steps than ContinuousRTC, especially under large delays with execution horizon as discussed in the Sec.~\ref{sec:method}.
Moreover, regarding the average solve rates of extended variants in Kinetix, we show three important obervations:
(1) First, DiscreteRTC can outperform \textit{\textbf{Training-time Continuous RTC}}~\citep{black2025training}, which is fine-tuned with an explicit objective and recipe to improve inpainting, further demonstrating the strong performance of the fine-tuning free DiscreteRTC.
(2) Next, by using fixed unmasking steps during inference, \textit{\textbf{DiscreteRTC + Fixed Steps}} baseline can achieve better performance by generating more fine-grained actions with the same computation budget as discussed in Sec.~\ref{sec:method_lower}.
(3) Moreover, \textit{\textbf{DiscreteVLASH}}, which integrates the VLASH~\citep{tang2025vlash} can achieve more stable performance across varying delays but at the cost of performance drop under small delays.
This demonstrates that discrete diffusion can be seamlessly combined with advanced inference-time methods for further gains.


\subsection{Real-World Results}
\label{sec:exp_real}
Next, we validate DiscreteRTC on a real robot platform, focusing on dynamic tasks where inference latency and control frequency have crucial impact on the closed-loop performance. Please refer to Appendix~\ref{app: realexp} for full deployment details and visualizations.

\textbf{Setup.}
% [SENT] hardware and data collection
We use a UR5e arm with a Robotiq gripper and a wrist-mounted RGB camera.
% [SENT] VLA backbone and action head
Developed with StarVLA~\citep{community2026starvla}, all policies share a Qwen2.5-VL-3B-Instruct~\citep{yang2025qwen3} VLM paired with a layerwise cross-attention DiT~\citep{peebles2023scalable} action head.
% [SENT] deployment details
The robot operates at 20\,Hz runing on a single RTX 4090.
Regarding the baselines, we mainly consider the Sync and RTC policy with flow-matching and discrete diffusion architecture to better demonstrate our main claims.



% [PARA] Task descriptions.
% [SENT] motivation: isolate reactiveness under inference delay
We design two dynamic tasks to directly stress the reactiveness, \textbf{Dynamic Pick} and \textbf{Dynamic Place}, where the robot is asked to pick a moving object or place the object on a moving platform. This requires the robot to keep re-estimating the object pose and to execute smoothly.



% [PARA] Real-world results.
% [SENT] results with success rate and inference time

\textbf{Results.}
% [SENT] single-para summary: sync fails, DiscreteRTC wins on success rate, inference time asymmetry
Table~\ref{tab:real_results} summarizes the real-world results.
Both Sync baselines completely fail ($0\%$), confirming that reactive asynchronous execution is indispensable in these regimes.
In terms of success rate, DiscreteRTC outperforms ContinuousRTC with the largest gap on the continuously moving Dynamic Pick target where the reactiveness matters more.
In terms of inference time, the two asynchronous methods move in opposite directions: ContinuousRTC inflates flow-matching cost from $151$ to $256$\,ms, a near-$1.7\times$ overhead from the $\Pi$GDM, whereas DiscreteRTC \emph{reduces} discrete diffusion cost from $303$ to $206$\,ms since fewer tokens need to be unmasked during inference. Consequently, although Discrete Sync alone is slower than Continuous Sync with naive bin-discretization and no accelerating techniques, DiscreteRTC ends up both faster and more accurate than Continuous RTC---demonstrating that the benefits of RTC and discrete diffusion compound rather than conflict.

\begin{table}[t]
\centering
\caption{Real-world results on the two dynamic tasks. We report task success rate over 20 trials per task and average inference time (ms) on a single RTX 4090.}
\label{tab:real_results}
\begin{tabular}{lccc}
\toprule
\multirow{2}{*}{\textbf{Method}} & \multicolumn{2}{c}{\textbf{Success Rate} $\uparrow$}  & \multirow{2}{*}{\textbf{Inference Time} $\downarrow$} \\
\cmidrule(lr){2-3}
& Dynamic Place & Dynamic Pick & \\
\midrule
Continuous Sync                   & 0\%           & 0\%            & 151ms \\
Discrete Sync                     & 0\%           & 0\%            & 303ms \\
\midrule
Continuous RTC                    & 90\%          & 45\%           & 256ms ($\sim1.7 \times$) \\
\textbf{Discrete RTC (Ours)}       & \textbf{100\%} & \textbf{95\%} & \textbf{206ms ($\sim0.7 \times$)} \\
\bottomrule

\end{tabular}
\vskip -0.1in
\end{table}





%===============================================================================

\section{Related Work}
\label{sec:related}

We briefly review the most relevant works in this section with an extended discussion in Appendix~\ref{app:related}.

\textbf{Efficient VLA via Discrete Diffusion.}
Prior efforts to accelerate VLAs span system-level optimization~\citep{ma2025running} and action tokenization~\citep{pertsch2025fast, liu2026oat, intelligence2025pi_}.
However, tokenization-based methods still rely on autoregressive decoding at inference time, limiting throughput.
Discrete diffusion~\citep{lou2023discrete} replaces autoregressive generation with parallel unmasking, and several recent works apply it to VLA action heads~\citep{liang2025discrete, wen2025llada, wen2025dvla, chen2025unified, ye2025dream, chen2026dfm, song2026fast}.
While these works demonstrate the efficiency of discrete diffusion, none investigate its structural advantage for asynchronous execution, which is the focus of ours.

\textbf{Asynchronous Execution Strategy.}
When inference latency exceeds the control period, the robot must act while computing~\citep{xiao2020thinking}.
With action chunking, strategies range from naive chunk switching and temporal ensembling~\citep{zhao2023learning} to rejection-based methods~\citep{liu2024bidirectional}.
Real-time chunking (RTC)~\citep{black2025real} poses chunk transitions as inpainting via $\Pi$GDM guidance~\citep{song2023pseudoinverse}, with follow-up works improving guidance design~\citep{yang2026abpolicy, lu2026faster} and learning inpainting at training time~\citep{black2025training, wang2026real, liu2026learning}.
However, all of these methods operate within the flow-matching framework, 
while DiscreteRTC is the first to exploit the native inpainting capability of discrete diffusion policies for natural asynchronous execution.




% % [PARA] Discrete diffusion for control is underexplored.
% \textbf{Efficient VLA via Discrete Diffusion.}
% % Many previous efficient VLA work focuses on system-level optimization
% To train and run VLAs efficiently, many prior efforts have focused on system-level optimization~\citep{kim2024openvla}, such as applying CUDA graph compilation and operator fusion to achieve real-time VLA inference on consumer GPUs~\citep{ma2025running}.
% % Another line fcous on the action tokenization, since it can accelerate VLA-pre-training but suffer from the auto-regressive style inference, so just use flow-matching head
% Action tokenization is another promising direction that leverages the fully optimized VLM pipeline by converting continuous actions into discrete tokens~\citep{intelligence2025pi_}.
% Beyond trivial bin-based discretization, works like FAST~\citep{pertsch2025fast} and OAT~\citep{liu2026oat} compress action chunks into compact discrete token sequences, significantly reducing the required generation length.
% However, these tokenization-based approaches still rely on autoregressive decoding at inference time, which inherently limits parallelism and throughput~\citep{intelligence2025pi_}; as a result, many of them fall back to a continuous flow-matching action head for deployment and are only used for efficient large-scale VLA pre-training.

% % Discrete Diffusion is proved to be efficient in LLM
% Discrete diffusion~\citep{lou2023discrete} offers a promising alternative by replacing autoregressive generation with parallel unmasking.
% Several recent works have begun exploring this direction:
% \citet{liang2025discrete} apply discrete diffusion to the VLA action head, generating action chunks via iterative re-masking and parallel decoding;
% LLaDA-VLA~\citep{wen2025llada} and DVLA~\citep{wen2025dvla} build unified discrete-diffusion VLA frameworks;
% UD-VLA~\citep{chen2025unified} unifies observation reasoning and action generation within a single discrete diffusion backbone.
% Among them, DreamVLA~\citep{ye2025dream} further replaces the VLM module with a discrete diffusion architecture, making it possible to pre-train the VLA with a unified backbone.


% % [PARA] Asynchronous execution methods.
% \textbf{Asynchronous Execution Strategy.}
% When inference latency exceeds the control period, the robot must act while the policy is still computing.
% The earliest formulation of this concurrent-control setting is Thinking While Moving~\citep{xiao2020thinking}, which executes the previous action during the current inference.
% With action chunking~\citep{lai2025action}, the simplest strategy is naive asynchronous execution, which switches to a new chunk as soon as it is ready, but this can cause inter-chunk discontinuity.
% Temporal ensembling~\citep{zhao2023learning} smooths consecutive chunks by averaging overlapping actions, at the cost of blurring distinct modes.
% Bidirectional decoding (BID)~\citep{liu2024bidirectional} uses rejection sampling to maintain cross-chunk continuity, but at substantially higher computational cost.

% Real-time chunking (RTC)~\citep{black2025real} elegantly poses chunk transitions as inpainting via $\Pi$GDM~\citep{song2023pseudoinverse} guidance, achieving strong performance on dynamic tasks.
% Following RTC, several works have focused on improving inference-time inpainting through better guidance design:
% ABPolicy~\citep{yang2026abpolicy} formulates the inpainting problem in a B-spline control-point action space for smooth manipulation.
% Orthogonally, another line of follow-up works focuses on training-time methods that directly learn to inpaint during fine-tuning:
% Training-Time RTC~\citep{black2025training} conditions on action prefixes during training to internalize the inpainting capability;
% REMAC~\citep{wang2026real} follows a similar strategy but further eases the training--inference mismatch by mixing ground-truth and predicted actions during fine-tuning;
% Legato~\citep{liu2026learning} learns a continuation-aware flow that reshapes the denoising dynamics to account for known prefixes at training time.
% FASTER~\citep{lu2026faster} adopts a similar idea to rethink the indivisible characteristic of flow-matching training and inference, introducing a horizon-aware schedule to accelerate the time for the first available action.
% There are also works that bridge both stages: VLASH~\citep{tang2025vlash} achieves lower policy-input temporal mismatch by estimating future robot states and pre-training the VLA to predict action chunks conditioned on both current observations and future states.

% Our method, DiscreteRTC, is the first effort to step outside the flow-matching policy framework and analyze the asynchronous execution problem for discrete diffusion policies. It requires no training-time modification, yet achieves strong performance by exploiting the native inpainting capability of discrete diffusion---and can be combined with methods like VLASH~\citep{tang2025vlash} for further gains.




%===============================================================================

\section{Discussion and Future Steps}
\label{sec:conclusion}
% [PARA] Conclusions: summarize the contribution and key results.
% [SENT] DiscreteRTC exploits native inpainting of discrete diffusion for async RTC
\textbf{Conclusions.}
We presented DiscreteRTC, which exploits the native inpainting capability of discrete diffusion for asynchronous real-time control, eliminating the need for inpainting-specific fine-tuning, heuristic guidance weights, and extra inference cost inherent to flow-matching-based RTC.
% [SENT] experimental validation and compatibility with VLASH
Simulated and real-world experiments show that DiscreteRTC outperforms ContinuousRTC and other asynchronous baselines, and can be combined with methods such as VLASH for further gains.

% [PARA] Limitations framed as shared open challenges that the community is actively addressing.
\textbf{Limitations.}
We note several limitations of our current instantiation, each of which coincides with an active research direction whose progress will directly benefit DiscreteRTC.
(1)~Our naive $k$-bin quantization ignores temporal structure and yields overly long token sequences, motivating compact, temporally aware tokenizers~\citep{pertsch2025fast, liu2026oat}.
(2)~Our modularized AR-VLM + discrete diffusion head prevents the backbone from participating in iterative unmasking, a gap that unified discrete-diffusion VLAs~\citep{wen2025llada, wen2025dvla, chen2025unified, ye2025dream} aim to close.
(3)~Traditional max-confidence unmasking does not yet translate the natural schedule into consistent gains.

\textbf{Future Steps.}
These limitations point to three future research directions: 1) \emph{A time-causally ordered action tokenizer} that produces temporally ordered, compact token representations; 2) \emph{A unified discrete diffusion VLA} in which observation reasoning and action generation share the same backbone; and 3) \emph{An appropriate yet principled unmasking strategy} that fully unleash the potential of natural schedule. Emerging techniques such as AR-block decoding align better with the implicit autoregressive structure induced by RTC and integrate seamlessly.




%===============================================================================

% \clearpage
% The acknowledgments are automatically included only in the final and preprint versions of the paper.
% \acknowledgments{If a paper is accepted, the final camera-ready version will (and probably should) include acknowledgments. All acknowledgments go at the end of the paper, including thanks to reviewers who gave useful comments, to colleagues who contributed to the ideas, and to funding agencies and corporate sponsors that provided financial support.}

%===============================================================================

% no \bibliographystyle is required, since the corl style is automatically used.
\bibliography{example}  % .bib



\newpage
\appendix
%===============================================================================
\section{Extended Related Works}
\label{app:related}

% [PARA] Discrete diffusion for control is underexplored.
\textbf{Efficient VLA via Discrete Diffusion.}
% Many previous efficient VLA work focuses on system-level optimization
To train and run VLAs efficiently, many prior efforts have focused on system-level optimization~\citep{kim2024openvla}, such as applying CUDA graph compilation and operator fusion to achieve real-time VLA inference on consumer GPUs~\citep{ma2025running}.
% Another line fcous on the action tokenization, since it can accelerate VLA-pre-training but suffer from the auto-regressive style inference, so just use flow-matching head
Action tokenization is another promising direction that leverages the fully optimized VLM pipeline by converting continuous actions into discrete tokens~\citep{intelligence2025pi_}.
Beyond trivial bin-based discretization, works like FAST~\citep{pertsch2025fast} and OAT~\citep{liu2026oat} compress action chunks into compact discrete token sequences, significantly reducing the required generation length.
However, these tokenization-based approaches still rely on autoregressive decoding at inference time, which inherently limits parallelism and throughput~\citep{intelligence2025pi_}; as a result, many of them fall back to a continuous flow-matching action head for deployment and are only used for efficient large-scale VLA pre-training.

% Discrete Diffusion is proved to be efficient in LLM
Discrete diffusion~\citep{lou2023discrete} offers a promising alternative by replacing autoregressive generation with parallel unmasking.
Several recent works have begun exploring this direction:
\citet{liang2025discrete} apply discrete diffusion to the VLA action head, generating action chunks via iterative re-masking and parallel decoding;
LLaDA-VLA~\citep{wen2025llada} and DVLA~\citep{wen2025dvla} build unified discrete-diffusion VLA frameworks;
UD-VLA~\citep{chen2025unified} unifies observation reasoning and action generation within a single discrete diffusion backbone.
Among them, DreamVLA~\citep{ye2025dream} further replaces the VLM module with a discrete diffusion architecture, making it possible to pre-train the VLA with a unified backbone.


% [PARA] Asynchronous execution methods.
\textbf{Asynchronous Execution Strategy.}
When inference latency exceeds the control period, the robot must act while the policy is still computing.
The earliest formulation of this concurrent-control setting is Thinking While Moving~\citep{xiao2020thinking}, which executes the previous action during the current inference.
With action chunking~\citep{lai2025action}, the simplest strategy is naive asynchronous execution, which switches to a new chunk as soon as it is ready, but this can cause inter-chunk discontinuity.
Temporal ensembling~\citep{zhao2023learning} smooths consecutive chunks by averaging overlapping actions, at the cost of blurring distinct modes.
Bidirectional decoding (BID)~\citep{liu2024bidirectional} uses rejection sampling to maintain cross-chunk continuity, but at substantially higher computational cost.

Real-time chunking (RTC)~\citep{black2025real} elegantly poses chunk transitions as inpainting via $\Pi$GDM~\citep{song2023pseudoinverse} guidance, achieving strong performance on dynamic tasks.
Following RTC, several works have focused on improving inference-time inpainting through better guidance design:
ABPolicy~\citep{yang2026abpolicy} formulates the inpainting problem in a B-spline control-point action space for smooth manipulation.
Orthogonally, another line of follow-up works focuses on training-time methods that directly learn to inpaint during fine-tuning:
Training-Time RTC~\citep{black2025training} conditions on action prefixes during training to internalize the inpainting capability;
REMAC~\citep{wang2026real} follows a similar strategy but further eases the training--inference mismatch by mixing ground-truth and predicted actions during fine-tuning;
Legato~\citep{liu2026learning} learns a continuation-aware flow that reshapes the denoising dynamics to account for known prefixes at training time.
FASTER~\citep{lu2026faster} adopts a similar idea to rethink the indivisible characteristic of flow-matching training and inference, introducing a horizon-aware schedule to accelerate the time for the first available action.
There are also works that bridge both stages: VLASH~\citep{tang2025vlash} achieves lower policy-input temporal mismatch by estimating future robot states and pre-training the VLA to predict action chunks conditioned on both current observations and future states.

Our method, DiscreteRTC, is the first effort to step outside the flow-matching policy framework and analyze the asynchronous execution problem for discrete diffusion policies. It requires no training-time modification, yet achieves strong performance by exploiting the native inpainting capability of discrete diffusion---and can be combined with methods like VLASH~\citep{tang2025vlash} for further gains.


%===============================================================================

\newpage
\section{Additional Experiments}
\label{app: addexp}
\subsection{Natural Schedule Ablations}
\label{app: naturalmask}

\begin{figure}[h]
    \centering
    \includegraphics[width=\linewidth]{figures/NaturalSchedule.png}
    \caption{\textbf{Unmasking Trajectory Sample with Natural Schedule Inference.} Green blocks denote unmasked action tokens, yellow blocks denote masked tokens, and the red rectangle marks the early-stop boundary beyond which tokens do not need to be unmasked before the next inference.}
    \label{fig:natural_schedule}
\end{figure}

In practice, the natural schedule does work as expected compared to the simple hard mask approach. In this section, we show how the inappropriate unmasking-schedule limits the performance of natural schedule.
We visualize the per-step unmasking trajectory of the discrete diffusion policy under the max-confidence decoding strategy~\citep{liang2025discrete}.
In the visualization, green blocks represent unmasked action tokens, yellow blocks represent masked tokens, and the red rectangle marks the early-stop boundary: only tokens inside this boundary need to be fully unmasked before the next inference, while tokens outside can remain partially masked and carry forward as the natural guidance signal.
For each generation, the policy starts with the first $d=4$ actions already unmasked (the frozen prefix from the previous chunk) and is expected to inpaint the remaining tokens in an implicitly autoregressive fashion---unmasking the early actions first so that execution can begin as soon as possible.

However, as shown in Figure~\ref{fig:natural_schedule}, the max-confidence unmasking trajectory does \emph{not} exhibit this desired implicit AR-style ordering.
Instead, the policy distributes its unmasking budget across the entire chunk based on per-token confidence, and in most cases still consumes the full 8 unmasking steps before all tokens inside the early-stop boundary are resolved.
This both neutralizes the inference-cost savings promised by early stopping and degrades follow-up generations: because more tokens are fully unmasked by the time inference completes, the next chunk inherits a larger committed prefix and correspondingly less flexibility from the natural guidance signal.
This observation motivates the third direction in our future steps---designing a principled unmasking strategy (\eg, AR-block decoding) that aligns with the implicit autoregressive structure induced by RTC and fully unleashes the potential of the natural schedule.


% Natural Performance
\newpage

\subsection{Detailed Main Results in Kinetix}
\begin{figure}[h]
    \centering
    \includegraphics[width=\linewidth]{figures/KinetixMainFull.pdf}
    \caption{\textbf{Detailed Main Results in Kinetix.}  The evaluation setup keeps the same with Figure~\ref{fig:kinetix_main}. The small plots represent the results per Kinetix environment}
    \label{fig:ft_ablation}
    \vskip -0.2in
\end{figure}


\subsection{Inpainting Fine-tuning}

\begin{figure}[h]
    \centering
    \includegraphics[width=\linewidth]{figures/FinetunedApp.pdf}
    \caption{\textbf{Fine-tuning Ablation in Kinetix.} The evaluation setup keeps the same with Figure~\ref{fig:kinetix_main}. The small plots represent the results per Kinetix environment}
    \label{fig:ft_ablation}
    \vskip -0.1in
\end{figure}

As an ablation of the fine-tuning-free claim, we explicitly fine-tune the pre-trained discrete diffusion policy with an inpainting-specific training recipe: for each training chunk we construct the target masking pattern by randomly masking as standard process first, then unmasking the prefix action tokens from $\{0, 1, 2, 3, 4\}$ to match the inference delay value distribution.
The intention is to bias the base policy toward the RTC-style prefix–continuation and further improve inpainting quality.


However, as results shown in Figure~\ref{fig:ft_ablation}, this naive fine-tuning strategy brings a consistent performance drop across all 12 Kinetix levels and every inference delay. This suggests that a naive prefix-unmasking recipe can actually \emph{hurts} the native inpainting capability of the pre-trained discrete diffusion policy. Improving beyond
the fine-tuning-free DiscreteRTC therefore appears to require either a carefully designed inpainting-aware fine-tuning recipe that does not distort the pre-training distribution, or just directly scaling the base policy's pre-training.

%===============================================================================
\newpage
\section{Kinetix Implementation Details}
\label{app: simexp}
The codebase is developed upon the \href{https://github.com/Physical-Intelligence/real-time-chunking-kinetix}{official RTC codebase}~\citep{black2025real}, implemented in JAX~\citep{jax2018github} with Flax NNx on 8 RTX 6000 PRO.

\textbf{Environment.}
We use Kinetix~\citep{matthews2024kinetix} (\texttt{Kinetix-Symbolic-Continuous-v1}, \texttt{large} configuration).
The physics simulation runs at 60\,Hz with frame skip 2, yielding 30\,Hz control.
Actions are normalized to $[-1, 1]$ with Gaussian noise (std 0.1) added during both data collection and evaluation.
Observations include a 4-frame history stack and the previous action.

\textbf{Policy Architecture.}
Both policies share an MLPMixer backbone with AdaLN conditioning: 4 blocks, each with token-mixing (hidden dim 64) and channel-mixing (hidden dim 512) paths, operating on $H{=}8$ tokens (one per chunk timestep) with channel dimension 256.

For \textbf{flow-matching policy}, the input is the observation concatenated with the noised action chunk.
A sinusoidal time embedding conditions the AdaLN layers, and the output is a velocity vector.
For \textbf{discrete diffusion policy}, actions are discretized into 512 bins via uniform quantization.
The input is the observation concatenated with bin embeddings of the partially masked token sequence.
The token sequence is packed from $(H, \text{action\_dim}, \text{embedding\_dim})$ to $(H, \text{channel\_dim})$ via a linear layer so the mixer sees the same 8-token structure.
A learned vector replaces the time conditioning.

\textbf{Training.}
Both policies are trained for 32 epochs with batch size 512 using AdamW.
The flow-matching policy uses a constant LR schedule and MSE loss on the predicted velocity field.
The discrete diffusion policy uses cosine LR decay and cross-entropy on masked positions plus an auxiliary L1 loss ($\lambda{=}0.1$) with a cosine masking schedule.

\textbf{Evaluation.}
Both policies use 5 denoising/unmasking steps.
For ContinuousRTC, $\Pi$GDM guidance with exponential-decay soft masking and clipping threshold $\beta$ steers generation toward the frozen prefix.
For DiscreteRTC, the frozen prefix tokens are directly placed, and for parallel evaluation across the environments, we use a hard mask for the chunk continuation.
All remaining hyperparameters are listed in Table~\ref{tab:hyperparams_sim}.

\begin{table}[h]
\centering
\caption{Simulation experiment hyperparameters.}
\label{tab:hyperparams_sim}
\begin{tabular}{lcc}
\toprule
\textbf{Hyperparameter} & \textbf{Flow-Matching} & \textbf{Discrete Diffusion} \\
\midrule
\multicolumn{3}{l}{\textit{Architecture}} \\
Channel / channel hidden / token hidden dim & \multicolumn{2}{c}{256 / 512 / 64} \\
Number of MLPMixer layers & \multicolumn{2}{c}{4} \\
Action chunk size ($H$) & \multicolumn{2}{c}{8} \\
Number of bins / embedding dim & --- & 512 / 128 \\
\midrule
\multicolumn{3}{l}{\textit{Training}} \\
Batch size / epochs & \multicolumn{2}{c}{512 / 32} \\
Optimizer / LR / weight decay & \multicolumn{2}{c}{AdamW / $3 \times 10^{-4}$ / $10^{-2}$} \\
Gradient norm clip & \multicolumn{2}{c}{10.0} \\
LR schedule (warmup steps) & Constant (1{,}000) & Cosine decay (2{,}000) \\
Training loss & MSE (velocity) & CE + L1 ($\lambda{=}0.1$) \\
Training mask schedule & --- & Cosine \\
\midrule
\multicolumn{3}{l}{\textit{Evaluation}} \\
Denoising / unmasking steps & 5 & 5 \\
Decode temperature / choice temperature & --- & 1.0 / 0.0 \\
RTC guidance & $\Pi$GDM ($\beta{=}5.0$) & Hard Mask \\
Inference delays ($d$) & \multicolumn{2}{c}{\{0, 1, 2, 3, 4\}} \\
Execution horizon ($s$) / trials & \multicolumn{2}{c}{$\max(1, d)$ / 2{,}048} \\
\bottomrule
\end{tabular}
\end{table}

%===============================================================================
\newpage
\section{Real-world Implementation Details}
\label{app: realexp}

\begin{figure}[h]
    \centering
    \includegraphics[width=0.9\linewidth]{figures/place.jpg}
    \caption{dynamic Pick and Place Real-world Setup}
    \label{fig:real_setup}
\end{figure}

\textbf{Hardware and Data.}
We use a single UR5e arm with a Robotiq gripper and a wrist-mounted RGB camera.
Demonstrations are recorded at 500\,Hz via the FastUMI pipeline.
Each action is a 10D vector $[\Delta x, \Delta y, \Delta z, \text{rot6d}(6), \text{gripper}]$, with translational dimensions normalized to $[-1, 1]$ via min-max scaling, rotation dimensions left unnormalized, and the gripper binarized to $\{0, 1\}$.
Images are center-cropped and resized to $224 \times 224$.
Training uses 8 RTX Pro 6000 with DeepSpeed ZeRO-2 and gradient checkpointing; deployment runs on a single RTX 4090.

\textbf{Policy Architecture.}
All policies use Qwen2.5-VL-3B-Instruct~\citep{yang2025qwen3} as the VLM backbone.
Both action heads share the same QwenPI-style layerwise cross-attention DiT~\citep{peebles2023scalable}: 36 blocks, each with AdaLN, cross-attention to the corresponding VLM hidden states, self-attention, and a feedforward network.
Both prepend 32 learnable future vision tokens and add positional embeddings.
When proprioceptive state is available, a two-layer MLP encodes it and prepends it to the sequence.

For \textbf{flow-matching head}, a three-layer MLP fuses per-timestep action embeddings with sinusoidal flow-time embeddings (output dim 2048); a sinusoidal \texttt{TimestepEncoder} provides AdaLN conditioning; and a two-layer MLP decodes the velocity output.
For \textbf{discrete diffusion head}, actions are quantized into 256 bins per dimension, yielding $16 \times 10 = 160$ tokens per chunk.
Tokens (including a learnable \texttt{[MASK]}) are embedded via \texttt{nn.Embedding}(257, 2048) with no timestep conditioning.

% \textbf{Training.}
% Both heads share the same optimizer (AdamW, $\beta_1{=}0.9$, $\beta_2{=}0.95$), learning rate schedule (cosine with warmup), and training budget (30k steps, effective batch size 64).
% The action head forward pass is repeated twice per VLM forward to amortize the VLM cost for both heads.
% The flow-matching head is trained with MSE on the predicted velocity field using a Beta($1.5, 1.0$) noise schedule.
% The discrete diffusion head is trained with cross-entropy on masked positions plus an auxiliary L1 loss ($\lambda{=}0.1$) using a cosine masking schedule.
% Mixed precision is used: the VLM runs in bfloat16, while both action heads run in float32.

\textbf{Deployment.}
The robot operates in closed-loop at 20\,Hz with 5$\times$ linear interpolation producing a 100\,Hz servo stream.
Camera capture and model inference run in a background thread, overlapping with action execution.
The rotation speed of the turntable for dynamic pick and place are 15 and 20 rpm.
The discrete diffusion head uses 8 MaskGIT unmasking steps with a cosine decode schedule (decode temperature 0.0, choice temperature 0.1).
By default we use non-fixed unmasking steps for the discrete diffusion policies and running RTC with $d=4$ and $s=4$.

All hyperparameters are summarized in Table~\ref{tab:real_hyperparams}.



\begin{table}[h]
\centering
\caption{Real-world experiment hyperparameters.}
\label{tab:real_hyperparams}
\begin{tabular}{lcc}
\toprule
\textbf{Hyperparameter} & \textbf{Flow-Matching} & \textbf{Discrete Diffusion} \\
\midrule
\multicolumn{3}{l}{\textit{VLM Backbone \& Action Head}} \\
VLM backbone & \multicolumn{2}{c}{Qwen2.5-VL-3B-Instruct} \\
Action head & \multicolumn{2}{c}{Layerwise Cross-Attention DiT} \\
DiT layers / heads / head dim & \multicolumn{2}{c}{36 / 32 / 64 (inner dim 2048)} \\
Dropout / future vision tokens & \multicolumn{2}{c}{0.2 / 32} \\
Action encoder & 3-layer MLP + sinusoidal time & Embedding(257, 2048) \\
Action decoder & 2-layer MLP (hidden 1024) & Linear $\rightarrow$ 256 logits \\
Timestep conditioning & TimestepEncoder (AdaLN) & None (zero vector) \\
\midrule
\multicolumn{3}{l}{\textit{Action Representation}} \\
Action dim / chunk length $H$ & \multicolumn{2}{c}{10 / 16} \\
Representation & Continuous & 256-bin quantization \\
\midrule
\multicolumn{3}{l}{\textit{Training}} \\
Optimizer & \multicolumn{2}{c}{AdamW ($\beta_1{=}0.9, \beta_2{=}0.95$)} \\
LR (VLM / action model) & \multicolumn{2}{c}{$10^{-5}$ / $10^{-4}$} \\
LR scheduler (warmup / total steps) & \multicolumn{2}{c}{Cosine, min $5{\times}10^{-7}$ (5k / 30k)} \\
Batch size (per device $\times$ GPUs) & \multicolumn{2}{c}{$8 \times 8 = 64$} \\
Mixed precision & \multicolumn{2}{c}{bfloat16 (VLM) / float32 (action head)} \\
Training loss & MSE (velocity) & CE + L1 ($\lambda{=}0.1$) \\
Noise / mask schedule & Beta($1.5, 1.0$), $s{=}0.999$ & Cosine \\
\midrule
\multicolumn{3}{l}{\textit{Inference \& Deployment}} \\
Denoising / unmasking steps & 8 & 8 \\
Decode / choice temperature & --- & 0.0 / 0.1 \\
Control / servo frequency & \multicolumn{2}{c}{20\,Hz / 100\,Hz ($5\times$ interpolation)} \\
Robot / camera & \multicolumn{2}{c}{UR5e + Robotiq / wrist-mounted fish-eye RGB} \\
\bottomrule
\end{tabular}
\end{table}






\end{document}
